% =====================================================================
\chapter{Variables aleatorias discretas}
\label{tema:5}
% =====================================================================

Una variable aleatoria asigna un número a cada resultado de un experimento; su distribución resume todo lo que interesa de ella. Este tema introduce la función de masa, la esperanza y la varianza, demuestra la linealidad de la esperanza y la explota mediante variables indicadoras, y estudia las distribuciones discretas básicas: Bernoulli, binomial, geométrica, Poisson y categórica. Se siguen los capítulos~3 y~4 de Blitzstein y Hwang~\cite{bh}, el capítulo~2 de Bertsekas y Tsitsiklis~\cite{bt} (clases~5 a~7 de MIT 6.041~\cite{mit6041f10}), el capítulo~4 de Ross~\cite{ross} y el capítulo~18, «Random Variables», de~\cite{mcs}.

\section{Variables aleatorias y función de masa}

\begin{definicion}[Variable aleatoria discreta]
Una \emph{variable aleatoria} es una función $X\colon\Omega\to\R$. Es \emph{discreta} si toma un conjunto finito o numerable de valores. Su \emph{función de masa} es $p_X(x)=\PP(X=x)$, donde $\{X=x\}$ abrevia el suceso $\{\omega\in\Omega: X(\omega)=x\}$, y su \emph{función de distribución} es $F_X(x)=\PP(X\le x)$ para $x\in\R$.
\end{definicion}

La función de masa cumple $p_X(x)\ge0$ y $\sum_xp_X(x)=1$, porque los sucesos $\{X=x\}$ forman una partición de $\Omega$. Para todo conjunto $S\subseteq\R$, $\PP(X\in S)=\sum_{x\in S}p_X(x)$. La función de distribución es no decreciente, tiende a 0 en $-\infty$ y a 1 en $+\infty$, es continua por la derecha y, en el caso discreto, es escalonada, con saltos de altura $p_X(x)$ en cada valor $x$.

\begin{ejemplo}\label{ej:t5-suma}
Sea $S$ la suma de dos dados. Contando en los 36 resultados equiprobables, $p_S(s)=\frac{6-\abs{s-7}}{36}$ para $s=2,\dots,12$; por ejemplo, $p_S(7)=6/36$ y $p_S(2)=1/36$. La suma de todos los valores es $(1+2+3+4+5+6+5+4+3+2+1)/36=1$.
\end{ejemplo}

\section{Esperanza y linealidad}

\begin{definicion}[Esperanza]
La \emph{esperanza} (o media) de una variable discreta $X$ es
\[
\E[X]=\sum_{x}x\,p_X(x),
\]
siempre que la serie converja absolutamente. Es el centro de masas de la distribución.
\end{definicion}

\begin{proposicion}[Esperanza de una función: LOTUS]\label{prop:t5-lotus}
Si $X$ es discreta y $g\colon\R\to\R$, entonces $\E[g(X)]=\sum_xg(x)\,p_X(x)$. Si $\Omega$ es discreto, además $\E[X]=\sum_{\omega\in\Omega}X(\omega)\,\PP(\{\omega\})$.
\end{proposicion}

\begin{proof}
Sea $Y=g(X)$. Para cada valor $y$, $\{Y=y\}$ es la unión disjunta de los $\{X=x\}$ con $g(x)=y$, luego $p_Y(y)=\sum_{x:g(x)=y}p_X(x)$ y
\[
\E[Y]=\sum_yy\sum_{x:g(x)=y}p_X(x)=\sum_y\sum_{x:g(x)=y}g(x)p_X(x)=\sum_xg(x)p_X(x),
\]
reordenando la serie, lo que es lícito por la convergencia absoluta. La segunda afirmación es el mismo argumento con $\Omega$ en el papel de los valores de $X$: $\{X=x\}$ es la unión disjunta de los $\{\omega\}$ con $X(\omega)=x$.
\end{proof}

\begin{teorema}[Linealidad de la esperanza]\label{teo:t5-linealidad}
Si $X$ e $Y$ son variables aleatorias discretas definidas sobre el mismo espacio y con esperanza finita, y $a,b\in\R$, entonces
\[
\E[aX+bY]=a\,\E[X]+b\,\E[Y].
\]
No se requiere ninguna hipótesis de independencia.
\end{teorema}

\begin{proof}
Se demuestra cuando $\Omega$ es discreto; el caso general es análogo usando la distribución conjunta. Por la \cref{prop:t5-lotus},
\[
\E[aX+bY]=\sum_{\omega}\bigl(aX(\omega)+bY(\omega)\bigr)\PP(\{\omega\})=a\sum_\omega X(\omega)\PP(\{\omega\})+b\sum_\omega Y(\omega)\PP(\{\omega\}),
\]
que es $a\E[X]+b\E[Y]$.
\end{proof}

\begin{definicion}[Variable indicadora]
La \emph{indicadora} de un suceso $A$ es la variable $I_A$ que vale 1 si ocurre $A$ y 0 si no.
\end{definicion}

\begin{proposicion}[Puente fundamental]\label{prop:t5-puente}
$\E[I_A]=\PP(A)$. En consecuencia, si $X=I_{A_1}+\dots+I_{A_n}$ cuenta cuántos de los sucesos $A_1,\dots,A_n$ ocurren, $\E[X]=\PP(A_1)+\dots+\PP(A_n)$.
\end{proposicion}

\begin{proof}
$\E[I_A]=1\cdot\PP(A)+0\cdot\PP(A^{c})=\PP(A)$. La segunda parte es la linealidad.
\end{proof}

\begin{ejemplo}[Puntos fijos]\label{ej:t5-fijos}
Se reparten al azar los sombreros de $n$ personas. Sea $X$ el número de personas que reciben el suyo. Con $A_i$ «la persona $i$ recibe su sombrero», $\PP(A_i)=(n-1)!/n!=1/n$, luego $\E[X]=n\cdot\frac1n=1$ para todo $n$. Calcular la distribución de $X$ es mucho más difícil (\cref{ej:t2-desarreglos}); la esperanza, gracias a las indicadoras, es inmediata, aunque los $A_i$ no son independientes~\cite[§4.4]{bh}.
\end{ejemplo}

\section{Varianza e independencia}

\begin{definicion}[Varianza]
La \emph{varianza} de $X$ es $\Var(X)=\E\bigl[(X-\mu)^{2}\bigr]$, con $\mu=\E[X]$, y su \emph{desviación típica} es $\sigma_X=\sqrt{\Var(X)}$. Mide la dispersión de $X$ alrededor de su media, en unidades al cuadrado.
\end{definicion}

\begin{proposicion}\label{prop:t5-var}
$\Var(X)=\E[X^{2}]-\E[X]^{2}$ y, para $a,b\in\R$, $\Var(aX+b)=a^{2}\Var(X)$. En particular $\Var(X)\ge0$ implica $\E[X^{2}]\ge\E[X]^{2}$.
\end{proposicion}

\begin{proof}
Por linealidad, $\E[(X-\mu)^{2}]=\E[X^{2}]-2\mu\E[X]+\mu^{2}=\E[X^{2}]-\mu^{2}$. Como $\E[aX+b]=a\mu+b$, se tiene $\Var(aX+b)=\E[(aX+b-a\mu-b)^{2}]=a^{2}\E[(X-\mu)^{2}]$.
\end{proof}

\begin{definicion}[Independencia de variables]
Las variables discretas $X_1,\dots,X_n$ son \emph{independientes} si $\PP(X_1=x_1,\dots,X_n=x_n)=\prod_i\PP(X_i=x_i)$ para todos los valores $x_1,\dots,x_n$.
\end{definicion}

\begin{proposicion}\label{prop:t5-producto}
Si $X$ e $Y$ son discretas, independientes y con esperanza finita, $\E[XY]=\E[X]\E[Y]$. En consecuencia, si $X_1,\dots,X_n$ son independientes (basta dos a dos), $\Var(X_1+\dots+X_n)=\Var(X_1)+\dots+\Var(X_n)$.
\end{proposicion}

\begin{proof}
Por LOTUS aplicada al par $(X,Y)$ y la independencia,
\[
\E[XY]=\sum_{x,y}xy\,\PP(X=x,Y=y)=\sum_x\sum_yxy\,p_X(x)p_Y(y)=\Bigl(\sum_xx\,p_X(x)\Bigr)\Bigl(\sum_yy\,p_Y(y)\Bigr).
\]
Para la varianza, sean $\mu_i=\E[X_i]$ y $S=\sum_iX_i$. Entonces $S-\E[S]=\sum_i(X_i-\mu_i)$ y
\[
\Var(S)=\E\Bigl[\sum_{i,j}(X_i-\mu_i)(X_j-\mu_j)\Bigr]=\sum_i\Var(X_i)+\sum_{i\neq j}\E[X_i-\mu_i]\,\E[X_j-\mu_j]=\sum_i\Var(X_i),
\]
porque para $i\neq j$ las variables $X_i-\mu_i$ y $X_j-\mu_j$ son independientes y tienen media~0.
\end{proof}

\section{Distribuciones discretas básicas}

\begin{definicion}[Bernoulli y binomial]
$X\sim\Ber(p)$, con $p\in[0,1]$, si $\PP(X=1)=p$ y $\PP(X=0)=1-p$: es la indicadora de un «éxito» de probabilidad $p$. $X\sim\Bin(n,p)$ si $X$ es el número de éxitos en $n$ ensayos de Bernoulli independientes con la misma probabilidad $p$.
\end{definicion}

\begin{teorema}\label{teo:t5-binomial}
Si $X\sim\Bin(n,p)$, entonces $\PP(X=k)=\binom nk p^{k}(1-p)^{n-k}$ para $k=0,\dots,n$, $\E[X]=np$ y $\Var(X)=np(1-p)$.
\end{teorema}

\begin{proof}
Cada sucesión concreta de resultados con $k$ éxitos y $n-k$ fracasos tiene probabilidad $p^{k}(1-p)^{n-k}$ por la independencia, y hay $\binom nk$ sucesiones así (se eligen las posiciones de los éxitos); por la aditividad se obtiene la función de masa, que suma~1 por el binomio de Newton. Para los momentos, $X=I_1+\dots+I_n$ con $I_j\sim\Ber(p)$ independientes, $\E[I_j]=p$ y $\Var(I_j)=\E[I_j^{2}]-p^{2}=p-p^{2}$, porque $I_j^{2}=I_j$. Se aplican la linealidad y la \cref{prop:t5-producto}.
\end{proof}

\begin{definicion}[Geométrica]
$X\sim\Geom(p)$, con $0<p\le1$, si $X$ es el número de ensayos de Bernoulli independientes de parámetro $p$ necesarios hasta obtener el primer éxito (incluido este). Su función de masa es $\PP(X=k)=(1-p)^{k-1}p$ para $k=1,2,\dots$
\end{definicion}

\begin{proposicion}\label{prop:t5-geom}
Si $X\sim\Geom(p)$: (i) $\PP(X>k)=(1-p)^{k}$; (ii) $\E[X]=1/p$ y $\Var(X)=(1-p)/p^{2}$; (iii) \emph{falta de memoria}: $\PP(X>m+k\mid X>m)=\PP(X>k)$ para enteros $m,k\ge0$.
\end{proposicion}

\begin{proof}
(i) $X>k$ si y solo si los $k$ primeros ensayos son fracasos. (ii) Con $q=1-p$, derivando la serie geométrica $\sum_{k\ge0}x^{k}=1/(1-x)$ para $\abs{x}<1$ se obtiene $\sum_{k\ge1}kx^{k-1}=1/(1-x)^{2}$, luego $\E[X]=p\sum_{k\ge1}kq^{k-1}=p/p^{2}=1/p$. Derivando otra vez, $\sum_{k\ge2}k(k-1)x^{k-2}=2/(1-x)^{3}$, de donde $\E[X(X-1)]=pq\cdot2/p^{3}=2q/p^{2}$ y $\Var(X)=\E[X(X-1)]+\E[X]-\E[X]^{2}=\frac{2q}{p^{2}}+\frac1p-\frac1{p^{2}}=\frac{q}{p^{2}}$. (iii) $\PP(X>m+k\mid X>m)=q^{m+k}/q^{m}=q^{k}$.
\end{proof}

\begin{definicion}[Poisson]
$X\sim\Poi(\lambda)$, con $\lambda>0$, si $\PP(X=k)=e^{-\lambda}\lambda^{k}/k!$ para $k=0,1,2,\dots$ Se cumple $\E[X]=\Var(X)=\lambda$ (\cref{eje:5-poisson}).
\end{definicion}

La distribución de Poisson modela el número de sucesos raros entre muchas oportunidades independientes: erratas en un libro, llegadas de peticiones a un servidor en un minuto, apariciones de una palabra poco frecuente en un texto largo. El resultado siguiente lo justifica.

\begin{teorema}[Aproximación de Poisson a la binomial]\label{teo:t5-poisson}
Sea $\lambda>0$ y $p_n=\lambda/n$. Para cada $k\ge0$ fijo,
\(
\lim_{n\to\infty}\binom nk p_n^{k}(1-p_n)^{n-k}=e^{-\lambda}\frac{\lambda^{k}}{k!}.
\)
\end{teorema}

\begin{proof}
Se escribe
\[
\binom nk\Bigl(\frac\lambda n\Bigr)^{k}\Bigl(1-\frac\lambda n\Bigr)^{n-k}=\frac{\lambda^{k}}{k!}\cdot\frac{n(n-1)\cdots(n-k+1)}{n^{k}}\cdot\Bigl(1-\frac\lambda n\Bigr)^{n}\cdot\Bigl(1-\frac\lambda n\Bigr)^{-k}.
\]
El segundo factor es $\prod_{j=0}^{k-1}(1-j/n)\to1$; el tercero tiende a $e^{-\lambda}$, por el límite $(1+x/n)^{n}\to e^{x}$; el cuarto tiende a~1 porque $k$ está fijo.
\end{proof}

\begin{definicion}[Categórica]
Sea $\V=\{v_1,\dots,v_K\}$ un conjunto finito y $\bm p=(p_1,\dots,p_K)$ con $p_i\ge0$ y $\sum_ip_i=1$. Una variable $X$ con valores en $\V$ tiene distribución \emph{categórica}, $X\sim\Cat(\bm p)$, si $\PP(X=v_i)=p_i$. Cuando los valores no son números, en lugar de la esperanza de $X$ se usan las de sus indicadoras: $\E[I_{\{X=v_i\}}]=p_i$.
\end{definicion}

\begin{table}[h]
\centering
\begin{tabular}{@{}lllll@{}}
\toprule
Distribución & Valores & $\PP(X=k)$ & $\E[X]$ & $\Var(X)$ \\
\midrule
$\Ber(p)$ & $\{0,1\}$ & $p^{k}(1-p)^{1-k}$ & $p$ & $p(1-p)$ \\
$\Bin(n,p)$ & $\{0,\dots,n\}$ & $\binom nk p^{k}(1-p)^{n-k}$ & $np$ & $np(1-p)$ \\
$\Geom(p)$ & $\{1,2,\dots\}$ & $(1-p)^{k-1}p$ & $1/p$ & $(1-p)/p^{2}$ \\
$\Poi(\lambda)$ & $\{0,1,\dots\}$ & $e^{-\lambda}\lambda^{k}/k!$ & $\lambda$ & $\lambda$ \\
Uniforme en $[n]$ & $\{1,\dots,n\}$ & $1/n$ & $(n+1)/2$ & $(n^{2}-1)/12$ \\
\bottomrule
\end{tabular}
\caption{Distribuciones discretas básicas.}
\label{tab:t5-discretas}
\end{table}

\begin{observacion}[Conexión con el curso]
La salida de un modelo de lenguaje es, en cada paso, una distribución categórica sobre el vocabulario, obtenida con la softmax (módulos~1 y~5); generar texto es muestrear de ella. La linealidad de la esperanza con indicadoras calcula, por ejemplo, el número esperado de tokens distintos en una muestra (\cref{eje:5-categorica}) o el número esperado de colisiones en una tabla \emph{hash} (\cref{eje:5-hash}, módulo~16). La métrica pass@$k$ del módulo~22 es una probabilidad calculada en un modelo hipergeométrico (\cref{eje:5-passk}).
\end{observacion}

\begin{observacion}[Errores frecuentes]
(1) Creer que la linealidad de la esperanza exige independencia: no la exige; la aditividad de la \emph{varianza} sí. (2) Escribir $\E[g(X)]=g(\E[X])$; en general es falso, por ejemplo $\E[X^{2}]\ge\E[X]^{2}$ con igualdad solo si $X$ es constante. (3) Escribir $\Var(aX)=a\Var(X)$ o $\Var(X-Y)=\Var(X)-\Var(Y)$: para $X$, $Y$ independientes, $\Var(X-Y)=\Var(X)+\Var(Y)$. (4) Confundir las dos versiones de la geométrica: la de este tema cuenta ensayos (valores $1,2,\dots$); otros textos cuentan fracasos (valores $0,1,\dots$), con media $(1-p)/p$. (5) Aplicar la binomial a extracciones \emph{sin} reemplazamiento, que no son independientes.
\end{observacion}

\begin{lecturas}
\item Blitzstein y Hwang~\cite{bh}, cap.~3, §3.1 a §3.3 y §3.8, y cap.~4, §4.1 a §4.4, §4.6 y §4.7.
\item Bertsekas y Tsitsiklis~\cite{bt}, §2.1 a §2.4 (clases~5 y~6 de MIT 6.041~\cite{mit6041f10}).
\item Ross~\cite{ross}, cap.~4, §4.1 a §4.7 y §4.9.
\item Lehman, Leighton y Meyer~\cite{mcs}, cap.~18, §18.1 a §18.5, en especial §18.5 «Linearity of Expectation».
\item Grinstead y Snell~\cite{gs}, §5.1 «Important Distributions» y §6.1 y §6.2.
\end{lecturas}

\section*{Ejercicios del tema 5}
\addcontentsline{toc}{section}{Ejercicios del tema 5}

\begin{ejercicio}{1}{eje:5-masa}
Sea $X$ con $\PP(X=k)=ck$ para $k=1,2,3,4$. Halle $c$, la función de distribución $F_X$ en los puntos $1,2,3,4$, $\E[X]$ y $\Var(X)$.
\end{ejercicio}

\begin{ejercicio}{2}{eje:5-diferencia}
Se lanzan dos dados equilibrados de tres caras (valores 1, 2 y 3) y se define $X$ como el valor del primero menos el del segundo. (a) Halle la función de masa de $X$, $\E[X]$ y $\Var(X)$. (b) Halle la función de masa de $X^{2}$ y compruebe que $\E[X^{2}]=\Var(X)$.
\fuente{adaptado de MIT 6.041~\cite{mit6041f10}, \emph{Problem Set}~3, problema~4.}
\end{ejercicio}

\begin{ejercicio}{2}{eje:5-binomial}
Un modelo resuelve cada uno de 10 problemas independientes con probabilidad $0{,}3$. Sea $X$ el número de problemas resueltos. Calcule $\PP(X=3)$, $\PP(X\ge1)$, $\E[X]$ y $\Var(X)$.
\end{ejercicio}

\begin{ejercicio}{2}{eje:5-geometrica}
Se repite un experimento con probabilidad de éxito $0{,}2$ hasta el primer éxito. Sea $X$ el número de intentos. Calcule $\PP(X=3)$, $\PP(X>5)$, $\E[X]$ y $\Var(X)$. Si ya se han fallado 3 intentos, ¿qué probabilidad hay de necesitar más de 8 en total?
\end{ejercicio}

\begin{ejercicio}{2}{eje:5-poisson-aprox}
En un texto de 1000 tokens, cada token es, independientemente, una palabra rara con probabilidad $0{,}002$. Sea $X$ el número de palabras raras. Calcule $\PP(X=0)$ y $\PP(X\ge2)$ con la binomial exacta y con la aproximación de Poisson, y compare.
\end{ejercicio}

\begin{ejercicio}{2}{eje:5-cumple-esp}
Hay 50 personas con cumpleaños uniformes e independientes entre 365 días. Use variables indicadoras para calcular (a) el número esperado de parejas de personas que cumplen años el mismo día; (b) el número esperado de días del año en que cumplen años al menos dos personas.
\fuente{adaptado de Harvard Stat 110~\cite{stat110}, \emph{Strategic Practice}~4, sección~2, problema~1.}
\end{ejercicio}

\begin{ejercicio}{2}{eje:5-hash}
Se insertan $k$ claves en una tabla \emph{hash} de $n$ casillas; cada clave va a una casilla uniforme, independientemente de las demás. Halle el número esperado de casillas vacías, de casillas con exactamente una clave y de casillas con más de una. Evalúe para $k=n=100$.
\fuente{adaptado de Harvard Stat 110~\cite{stat110}, \emph{Strategic Practice}~4, sección~2, problema~4.}
\end{ejercicio}

\begin{ejercicio}{2}{eje:5-hijos}
Una pareja tiene hijos hasta tener al menos un niño y una niña. Suponiendo sexos equiprobables e independientes, calcule el número esperado de hijos.
\fuente{adaptado de Harvard Stat 110~\cite{stat110}, \emph{Homework}~4, problema~3.}
\end{ejercicio}

\begin{ejercicio}{2}{eje:5-inversiones}
Una \emph{inversión} de una lista de números distintos es un par de posiciones $i<j$ con el elemento en $i$ mayor que el elemento en $j$. Si la lista de $n$ números distintos está en orden aleatorio uniforme, calcule el número esperado de inversiones. Compruébelo enumerando las $4!$ permutaciones de 4 elementos (o un subconjunto representativo).
\fuente{adaptado de Harvard Stat 110~\cite{stat110}, \emph{Homework}~4, problema~6(a).}
\end{ejercicio}

\begin{ejercicio}{2}{eje:5-categorica}
Un modelo genera 10 tokens independientes con distribución categórica $\PP(\texttt{a})=0{,}5$, $\PP(\texttt{b})=0{,}3$, $\PP(\texttt{c})=0{,}2$. (a) Calcule el número esperado de veces que aparece cada token y la varianza del número de apariciones de \texttt{a}. (b) Calcule el número esperado de tokens distintos que aparecen.
\end{ejercicio}

\begin{ejercicio}{2}{eje:5-passk}
Un modelo genera $n=10$ soluciones a un problema, de las que $c=3$ son correctas. Se eligen $k$ de las 10 al azar, sin reemplazamiento. (a) Demuestre que la probabilidad de que al menos una de las $k$ sea correcta es $1-\binom{n-c}{k}\big/\binom{n}{k}$. (b) Evalúe para $k=1,2,5$.
\end{ejercicio}

\begin{ejercicio}{3}{eje:5-colas}
(a) Demuestre que si $X$ toma valores en $\{0,1,2,\dots\}$, entonces $\E[X]=\sum_{k\ge1}\PP(X\ge k)$, escribiendo $X=\sum_{k\ge1}I_{\{X\ge k\}}$ o intercambiando el orden de una suma doble. (b) Úselo para obtener de nuevo la esperanza de una $\Geom(p)$.
\fuente{adaptado de MIT 6.041~\cite{mit6041f10}, \emph{Problem Set}~3, problema~3(a), y de Harvard Stat 110~\cite{stat110}, \emph{Strategic Practice}~4, sección~1, problema~5.}
\end{ejercicio}

\begin{ejercicio}{3}{eje:5-poisson}
Sea $X\sim\Poi(\lambda)$. Demuestre que su función de masa suma~1, que $\E[X]=\lambda$ y que $\E[X(X-1)]=\lambda^{2}$, y deduzca que $\Var(X)=\lambda$.
\end{ejercicio}

\begin{ejercicio}{3}{eje:5-sin-memoria}
Sea $X$ una variable con valores en $\{1,2,\dots\}$ y $\PP(X>1)\in(0,1)$, tal que $\PP(X>m+k\mid X>m)=\PP(X>k)$ para todos los enteros $m,k\ge0$ con $\PP(X>m)>0$. Demuestre que $X$ tiene distribución geométrica. \emph{Indicación:} estudie $g(k)=\PP(X>k)$.
\end{ejercicio}

\begin{ejercicio}{3}{eje:5-geom-cond}
Sean $X$ e $Y$ independientes con distribución $\Geom(p)$. Demuestre que, para $n\ge2$, la distribución de $X$ condicionada a $X+Y=n$ es uniforme en $\{1,\dots,n-1\}$.
\fuente{adaptado de MIT 6.041~\cite{mit6041f10}, \emph{Problem Set}~4, problema~3.}
\end{ejercicio}
